\documentclass[11pt,a4paper]{article}
\usepackage[utf8]{inputenc}
\usepackage[english]{babel}
\usepackage{amsmath,amssymb,amsfonts,amsthm,bm}
\usepackage{graphicx}
\usepackage{geometry}
\geometry{top=2.2cm, bottom=2.2cm, left=2.2cm, right=2.2cm}
\usepackage{hyperref}
\usepackage{booktabs}
\usepackage{cite}

\newtheorem{theorem}{Theorem}
\newtheorem{proposition}{Proposition}
\newtheorem{lemma}{Lemma}

\title{\textbf{Deductive Derivation of the Universal Feynman Path-Contour Master-Functional in Unified Phase Dynamics: From Dyson's Graviton Impossibility to Entire Schwinger Holonomies and Wheeler--DeWitt Closure}}

\author{\textbf{Hryhorii Pavliuk}\\[0.4em]
\small Independent Researcher \quad\textbar\quad \small ORCID: \href{https://orcid.org/0009-0000-0899-3463}{0009-0000-0899-3463}\\
\small Email: \texttt{hrpavliuk@gmail.com} \\[0.2em]
\small Monograph Concept DOI: \href{https://doi.org/10.5281/zenodo.24226129}{10.5281/zenodo.24226129} \quad Dedicated DOI: \href{https://doi.org/10.5281/zenodo.24386437}{10.5281/zenodo.24386437}}
\date{\today}

\begin{document}
\maketitle

\begin{abstract}
We present the complete deductive mathematical derivation of the universal path-contour master-functional:
\begin{equation}
\mathbf{\Psi}[\Sigma] = \frac{1}{\mathcal{N}} \mathrm{Tr}_{\mathcal{G}} \, \mathcal{P} \exp \left( \frac{i}{\hbar} \oint_{\partial \Sigma} \mathbf{\Pi}_\mu dx^\mu \right) \exp \left( - \frac{s_{\min}^2}{2\hbar^2} \int_\Sigma \mathcal{F}_{\mu\nu}^a \mathcal{F}^{\mu\nu}_a d\sigma \right), \quad \hat{\mathcal{H}} \, \mathbf{\Psi}[\Sigma] \equiv 0
\end{equation}
which serves as the sole axiomatic foundation of the \textit{Unified Phase Dynamics of Feynman Path Contours} (UPD). This master-functional is rigorously derived as the unique mathematically consistent resolution to the foundational crisis of theoretical physics: the non-renormalizability of perturbative quantum gravity and the unobservability of point-like gravitons (Freeman Dyson, 2013).

The deductive derivation proceeds through seven consecutive mathematical stages:
(1) The Perturbative Graviton Crisis and Dyson's Theorem; 
(2) Universal Phase Holonomies on Fiber Bundles; 
(3) The Unified Canonical Operator $\mathbf{\Pi}_\mu$ on $\mathcal{G} = U(1)_Y \times SU(2)_L \times SU(3)_c \times \mathrm{Spin}(3,1)$; 
(4) Non-Abelian Stokes Surface Transformation; 
(5) Entire Analytic Schwinger Regularization cancelling the 2-loop Goroff--Sagnotti divergence identically ($\mathrm{Res}_{\epsilon \to 0} \mathcal{A}_{\text{GS}}^{(2)} \equiv 0$) without Ostrogradsky ghosts; 
(6) Relational Quantum Time and Wheeler--DeWitt Constraint Closure; 
(7) Independence from Newton's Gravitational Constant $G$, demonstrating that $G = 12\pi c^3 s_{\min} / \hbar$ is an emergent macroscopic compliance coefficient of the vacuum phase substrate.
\end{abstract}

\section{Introduction: The Graviton Non-Renormalizability Crisis}
Standard quantum gravity expanding $g_{\mu\nu} = \eta_{\mu\nu} + \kappa h_{\mu\nu}$ possesses a coupling $[\kappa] = [M]^{-1} < 0$. Goroff and Sagnotti (1985) established that at two loops, pure Einstein gravity exhibits a non-renormalizable ultraviolet divergence:
\begin{equation}
\Delta \mathcal{L}^{(2)} = \frac{209}{2880 (4\pi)^4 \epsilon} \sqrt{-g} \, C_{\mu\nu\rho\sigma} C^{\rho\sigma\alpha\beta} C_{\alpha\beta}{}^{\mu\nu}
\end{equation}
Higher-derivative counterterms produce Ostrogradsky ghosts, destroying unitarity.

\section{Dyson's Theorem and Operational Graviton Invisibility}
\begin{theorem}[Freeman Dyson, 2013]
Any detector capable of registering a single graviton of wavelength $\lambda_g$ with unit probability must satisfy $R \le \lambda_g$ and $\lambda_g \le 4GM/c^2$, enforcing $R \le 4GM/c^2 = 2 r_s$. Thus, the detector necessarily collapses into a black hole before detection occurs.
\end{theorem}
Point-like gravitons are unobservable. Ontology must reside in non-local closed 1D path holonomies and 2D surfaces.

\section{Universal Phase Holonomies on Principal Fiber Bundles}
Open quantum path amplitudes transform non-locally under gauge transformations: $U[\gamma] \to U(x_2) U[\gamma] U^\dagger(x_1)$. Gauge invariance demands closed paths $\gamma(0) = \gamma(1)$:
\begin{equation}
W(\gamma) = \frac{1}{\mathcal{N}} \mathrm{Tr}_{\mathcal{G}} \left[ \mathcal{P} \exp \left( \frac{i}{\hbar} \oint_\gamma \mathbf{\Pi}_\mu dx^\mu \right) \right]
\end{equation}
The unified canonical operator on $\mathbf{V} = \mathbb{C}^1 \oplus \mathbb{C}^2 \oplus \mathbb{C}^3 \oplus \mathbb{C}^4$ is:
\begin{equation}
\mathbf{\Pi}_\mu = p_\mu \mathbf{I}_{10} + g_1 \mathcal{A}_\mu^{(Y)} \oplus g_2 \mathcal{A}_\mu^{(L)a}\frac{\tau_a}{2} \oplus g_3 \mathcal{A}_\mu^{(c)b}\frac{\lambda_b}{2} \oplus \hbar \omega_\mu^{(ab)}\Sigma_{ab}
\end{equation}

\section{Non-Abelian Stokes Theorem and Surface Holonomies}
\begin{equation}
\mathcal{P} \exp \left( \frac{i}{\hbar} \oint_{\partial \Sigma} \mathbf{\Pi}_\mu dx^\mu \right) = \mathcal{P}_\sigma \exp \left( \frac{i}{2\hbar} \int_\Sigma \mathcal{F}_{\mu\nu} d\sigma^{\mu\nu} \right)
\end{equation}
where $\mathcal{F}_{\mu\nu} = \partial_\mu \mathbf{\Pi}_\nu - \partial_\nu \mathbf{\Pi}_\mu + \frac{i}{\hbar}[\mathbf{\Pi}_\mu, \mathbf{\Pi}_\nu]$.

\section{Entire Analytic Schwinger Regulator}
To regularize loop divergences without Ostrogradsky ghosts, we introduce the entire form-factor:
\begin{equation}
\mathcal{E}_{s_{\min}}(\Box) = \exp \left( - \frac{s_{\min}^2 \Box^2}{\hbar^2} \right) \implies \exp \left( - \frac{s_{\min}^2}{2\hbar^2} \int_\Sigma \mathcal{F}_{\mu\nu}^a \mathcal{F}^{\mu\nu}_a d\sigma \right)
\end{equation}
Because $\mathcal{E}_{s_{\min}}(z)$ has no finite zeroes on $\mathbb{C}$, no ghost poles exist, reflection positivity (OS2) is preserved, and:
\begin{equation}
\mathrm{Res}_{\epsilon \to 0} \mathcal{A}_{\text{GS}}^{(2)} \equiv 0
\end{equation}

\section{Wheeler--DeWitt Constraint Closure and Emergent $G$}
The closed universe enforces $\hat{\mathcal{H}} \mathbf{\Psi}[\Sigma] \equiv 0$. The Dirac constraint algebra closes without anomalies: $[\hat{\mathcal{H}}(N_1), \hat{\mathcal{H}}(N_2)] \mathbf{\Psi} \equiv 0$.
The gravitational constant $G$ does not enter $\mathbf{\Psi}[\Sigma]$; it emerges as the macroscopic compliance of the vacuum phase substrate:
\begin{equation}
G = \frac{12\pi c^3 s_{\min}}{\hbar}
\end{equation}
In the flat-space limit of $SU(3)$ Yang--Mills, $\lim_{G \to 0} \Delta_{0^{++}} \approx 1.59\text{ GeV} > 0$.

\begin{thebibliography}{99}
\bibitem{PavliukMonograph} H.~Pavliuk, \textit{Unified Phase Dynamics of Feynman Path Contours}, Monograph, 2nd Ext. Ed., Zenodo Record 24386437, DOI: \href{https://doi.org/10.5281/zenodo.24386437}{10.5281/zenodo.24386437} (2026).
\bibitem{Dyson2013} F.~Dyson, Int. J. Mod. Phys. A \textbf{28}, 1330041 (2013).
\bibitem{GoroffSagnotti1985} M.~H.~Goroff and A.~Sagnotti, Phys. Lett. B \textbf{160}, 81--86 (1985).
\bibitem{Feynman1948} R.~P.~Feynman, Rev. Mod. Phys. \textbf{20}, 367--387 (1948).
\bibitem{Dirac1931} P.~A.~M.~Dirac, Proc. R. Soc. Lond. A \textbf{133}, 60--72 (1931).
\end{thebibliography}

\end{document}
